\documentclass[11pt,a4paper]{article}

\usepackage[utf8]{inputenc}
\usepackage[T1]{fontenc}
\usepackage[french,english]{babel}

\usepackage{amsmath,amssymb,amsthm}
\usepackage{mathrsfs}
\usepackage{graphicx}
\usepackage{hyperref}
\usepackage{geometry}
\usepackage{physics}
\usepackage{enumitem}
\setlist[itemize,1]{label=---}
\geometry{margin=2.5cm}
\usepackage{csquotes}

\newtheorem{theorem}{Theorem}[section]
\newtheorem{proposition}[theorem]{Proposition}
\newtheorem{lemma}[theorem]{Lemma}
\newtheorem{corollary}[theorem]{Corollary}

\theoremstyle{definition}
\newtheorem{definition}[theorem]{Definition}
\newtheorem{remark}[theorem]{Remark}
\newtheorem{example}[theorem]{Example}

\newcommand{\PDL}{\mathrm{PDL}}
\newcommand{\LDP}{\mathrm{LDP}}
\newcommand{\Gprot}{G_{\mathrm{p}}}
\newcommand{\etaL}{\eta} % leakage parameter
\newcommand{\Fopt}{\mathcal{F}} % optimisation functional
\newcommand{\Cconf}{\mathcal{C}} % configuration space
\newcommand{\Ttri}{\mathcal{T}} % set of triangles
\newcommand{\Prob}{\mathbb{P}}
\newcommand{\E}{\mathbb{E}}

\begin{document}

\title{Logical Leakage as Self-Maintained Probability:\\
A Topological Reformulation in the Projective Dynamic Logo Framework}

\author{C.~Laubscher\thanks{Independent researcher in theoretical and mathematical physics, Switzerland.}\\[4pt]}
\date{\small \today}

\maketitle

\begin{abstract}
The Projective Dynamic Logo (PDL/LDP) framework represents physical reality as a network of minimal logical closures on finite signed graphs, deriving particles and coupling constants from coherence constraints rather than from a pre-specified spacetime or fundamental fields. Within this framework, a small but irreducible leakage parameter \(\etaL\) arises at the proton level as the minimal fraction of frustrated (violated) triangles compatible with a hierarchical organisation into valence cores, relational sea, and active surface, and reappears in relational expressions for gravitational and thermal couplings. 

In this note, we show that the leakage functional \(\eta(G)\) can be reinterpreted as an intrinsic non-closure probability: for any finite signed graph \(G\), \(\eta(G)\) coincides with the probability, under a natural uniform measure on vertex triplets, that a randomly selected triangle is violated. At the proton level, \(\etaL\) thus acquires the status of a self-maintained probability, determined by the same coherence-optimisation principle that attempts to minimise it, rather than by an externally imposed stochastic postulate. 

On this basis, we introduce a topology of coexistence on the space of closed configurations, generated by logical compatibility relations, and define a coherence-cost pseudometric that quantifies the minimal additional leakage required to sustain multiple closures within a common neighbourhood. In this setting, the emergent metric previously proposed in PDL/LDP is reformulated as a cost-based distance, and gravitation is interpreted as a large-scale modulation of coherence cost induced by hierarchically filtered leakage. 

We then distinguish a local logical entropy, associated with the leakage probability within individual configurations, from a global topological entropy, associated with the distribution of effective leakage levels over the configuration space when many structures coexist. This yields a logical entropic paradox: although the optimisation functional drives each structure towards minimal leakage and hence low local entropy, the aggregation of many such nearly optimal closures into rich hierarchical assemblies generically gives rise to high global entropic complexity. 

Finally, we outline how this probabilistic and topological reformulation interfaces with existing PDL/LDP reinterpretations of Born-type probabilities, gravitational coupling, and Boltzmann’s constant, thereby suggesting that quantum statistics, gravitation, and thermodynamic granularity may be construed as distinct macroscopic manifestations of a single, logically enforced imperfection of closure in a finite relational network.
\end{abstract}

\newpage

\section{Introduction}

The Projective Dynamic Logo (PDL), or \textit{Logo Dynamique Projectif} (LDP), framework aims to reconstruct physical reality from purely relational axioms formulated on finite signed graphs, without positing \textit{a priori} a background spacetime manifold, predefined fields, or primitive material substrates. Within this approach, physical entities are represented as minimal logical closures: finite configurations of pulsating binary relations whose internal consistency is maintained by an optimisation principle rather than by externally prescribed dynamical laws. 

Subject to four fundamental axioms—minimal binary pulsation, triangular coherence, minimal completeness, and logical optimisation—the simplest admissible elementary stationary structure is given by the complete signed graph on four vertices and six edges, denoted \((4,6)\). This configuration supports a global two-step cycle of sign inversion while preserving exact triangular coherence throughout. It can thus be naturally interpreted as a logical prototype of an electron at rest. 

More complex composite structures, such as the proton, are modelled as hierarchical architectures constructed from such elementary closures. These architectures integrate locally stationary domains (valence cores), an extended relational sea, and a finite active surface that mediates couplings to electron-type \((4,6)\) blocks.

Within this framework, several constants that are conventionally regarded as fundamental empirical parameters are assigned a relational interpretation. The reduced Planck constant \(\hbar\) is construed as the minimal quantum of action associated with a single coherence cycle of an elementary closure; the fine-structure constant \(\alpha\) is interpreted as a ratio of internal to surface coherence on the proton’s active boundary; Boltzmann’s constant \(k_B\) is understood as a measure of the granularity of coherence within ensembles of hierarchically organised structures; and an effective gravitational coupling \(G_{\mathrm{eff}}\) is attributed to a minimal, hierarchically amplified leakage of coherence into an emergent metric background. 

A central role in these reinterpretations is played by a leakage parameter \(\etaL\), defined as the minimal fraction of violated triangles that persists in proton-type architectures even after coherence has been maximised under the relevant combinatorial constraints. This leakage is small but strictly positive, and it reappears—frequently in exponentiated form—in expressions for both gravitational and thermal coupling strengths.

The present note examines the logical interpretation of the leakage parameter and its connection to probability theory, topology, and metric structure. Instead of treating \(\etaL\) as a small ad hoc correction term or as an externally prescribed noise amplitude, we demonstrate that it can be reinterpreted as an intrinsic probability of non-closure: for a finite signed graph, the proportion of violated triangles coincides exactly with the probability, under a natural uniform measure, that a randomly selected triangle fails to close. At the proton level, the leakage parameter thus becomes a \emph{self-maintained probability}, determined by the same optimisation principle that drives its minimisation, and fully specified by the discrete organisation of valence blocks, relational sea, and active surface. From this standpoint, probability does not appear as an independent primitive concept; rather, it arises as the logically inevitable remnant of imperfect closure in finite relational structures.

On this probabilistic foundation, we explicitly construct a topology of coexistence on the space of closed configurations. Two structures are defined to be compatible if they can be jointly realized in a single configuration without increasing leakage beyond a prescribed admissible threshold. The corresponding compatibility neighbourhoods induce a topology, and a coherence-cost pseudometric is introduced by quantifying the minimal additional leakage necessary to retain configurations within a common neighbourhood. In this formulation, the emergent metric proposed in earlier PDL/LDP studies can be reformulated as a cost-based distance on a purely relational configuration space: what is conventionally interpreted as spatial separation is reinterpreted as a quantitative measure of differential coherence cost, while gravitation manifests as a large-scale modulation of this cost, governed by the hierarchical propagation of leakage.

The probabilistic and topological reformulation additionally enables a clear distinction between a local logical entropy, associated with the leakage probability within an individual configuration, and a global or topological entropy, associated with the distribution of effective leakage values over the entire configuration space when multiple structures coexist. At the local level, the optimisation functional minimises leakage and thereby reduces local entropy: the elementary \((4,6)\) closure exhibits zero leakage, and proton-type structures display only negligible leakage. At the global level, however, the aggregation of numerous such nearly optimal closures into complex hierarchies generates a broad spectrum of effective leakage magnitudes and, consequently, an increase in global entropic complexity. This tension—between the local tendency towards perfect closure and the global proliferation of distinct configurations that accommodate a small but irreducible degree of non-closure—will be referred to as a logical entropic paradox.

Finally, we outline how this reformulation interfaces with the physical quantities already analyzed within the PDL/LDP framework. The same logically self-sustaining leakage mechanism that underpins gravitational coupling and thermal granularity can be interpreted as the origin of Born-type probabilities, through its redistribution across mixed-triangle channels in electron–proton couplings. Although the present note does not provide a complete mathematical derivation of the Born rule, of the exponent \(18\) appearing in the gravitational coupling, or of the numerical value of \(\etaL\) from first principles, it advances a coherent conceptual schema within which these problems can be formulated with greater rigor. The overarching objective is not to introduce additional ad hoc assumptions, but rather to elucidate how coherence, closure, leakage, probability, topology, and metric structure are integrated within a unified relational framework, and to delineate concrete mathematical questions whose resolution would enhance the predictive capabilities of the PDL/LDP programme.

\section{Logical Preliminaries}

In this section, we summarise the minimal logical components of the PDL/LDP framework that are necessary for the subsequent probabilistic and topological developments. We do not aim to reproduce the complete derivation of the \((4,6)\) structure or of the proton hierarchy. Instead, we establish the relevant notation and highlight the foundational role played by triangular coherence, closure properties, and optimisation principles.

\subsection{Signed graphs and triangular coherence}

The fundamental objects of the framework are finite signed graphs, which are used to model networks of pairwise agreement/disagreement relations between informational entities.

\begin{definition}[Signed graph]
A \emph{signe\discretionary{d}{}{} graph} is a triple \(G = (V,E,\sigma)\) where:
\begin{itemize}
\item \(V\) is a finite set of vertices, representing informational entities;
\item \(E \subseteq \{ \{i,j\} : i,j \in V, i \neq j \}\) is a set of undirected edges;
\item \(\sigma : E \to \{+1,-1\}\) is a \emph{sign function} that assigns to each edge a sign, interpreted as a state of agreement \((+1)\) or disagreement \((-1)\).
\end{itemize}
We write \(|V|\) for the number of vertices and \(|E|\) for the number of edges.
\end{definition}

Within this framework, triangles—cycles of length three—play a distinguished role as minimal units of logical closure.

\begin{definition}[Triangles and coherence]
Given a finite signed graph \(G = (V,E,\sigma)\), the set of unordered vertex triplets (potential triangles) is
\[
 \Ttri(G) := \{ \{i,j,k\} \subset V : i,j,k \text{ distinct} \},
\]
so that \(|\Ttri(G)| = \binom{|V|}{3}\) whenever \(|V| \geq 3\). A triplet \(T = \{i,j,k\} \in \Ttri(G)\) is called:
\begin{itemize}
\item \emph{present} if all three edges \(\{i,j\}\), \(\{j,k\}\), and \(\{i,k\}\) belong to \(E\);
\item \emph{coherent} if it is present and the product of the signs on its edges satisfies
\[
 \sigma(\{i,j\})\,\sigma(\{j,k\})\,\sigma(\{i,k\}) = +1;
\]
\item \emph{violated} if it is not coherent, i.e., if at least one of its three edges is missing from \(E\), or if all three edges are present but the product of edge signs equals \(-1\).
\end{itemize}
\end{definition}

Coherent triangles thus constitute minimal closed loops of reference in which the pattern of agreement and disagreement can be traversed cyclically without generating an uncompensated inconsistency. By contrast, violated triangles indicate either a lack of relational support (missing edge) or an incompatible configuration of edge signs that prevents the establishment of full closure.

\begin{definition}[Fraction of violated triangles]
Let \(G = (V,E,\sigma)\) be a signed graph with \(|V| \geq 3\). The \emph{fraction of violated triangles} (also referred to as the leakage functional) \(\eta(G)\) is defined as
\[
 \eta(G) := \frac{\#\{ T \in \Ttri(G) : T \text{ is violated} \}}{|\Ttri(G)|} = \frac{\#\{ T \in \Ttri(G) : T \text{ is violated} \}}{\binom{|V|}{3}}.
\]
For \(|V| < 3\), we set \(\eta(G) := 0\) by convention.
\end{definition}

By definition, \(\eta(G) \in [0,1]\), with \(\eta(G) = 0\) if and only if every triangle in \(G\) is coherent and present. Signed graphs with small values of \(\eta(G)\) exhibit a high degree of triangular coherence and are, in this sense, natural candidates for stable closed configurations.

\subsection{Elementary closures and hierarchical architectures}

The PDL/LDP axioms entail that elementary stationary configurations must be closed, non-hierarchical, and internally symmetric, while maximising the density of coherent triangles. Under these constraints, one can show that the first admissible elementary closure is realised by a complete signed graph on four vertices, with six edges and an appropriate assignment of signs.

\begin{remark}[The \texorpdfstring{$(4,6)$}{(4,6)} structure]
Let \(K_4\) denote the complete graph on four vertices. There exists a sign assignment \(\sigma\) on the six edges of \(K_4\) such that every triangular subgraph is coherent and such that a global sign inversion \(\sigma \mapsto -\sigma\) preserves this triangular coherence. The resulting signed graph, denoted \(G_{4,6}\), thus satisfies \(|V| = 4\), \(|E| = 6\), and \(\eta(G_{4,6}) = 0\). Within the PDL/LDP framework, \(G_{4,6}\) realises the first non-trivial elementary stationary configuration and is interpreted as corresponding to an electron at rest.
\end{remark}

Composite structures, such as nucleons, are described as hierarchical assemblies of such elementary closures together with additional relational components. In particular, the proton is modelled by a two-level architecture that combines locally stationary domains with a dynamical relational background.

\begin{remark}[Proton hierarchy and leakage]
In the PDL/LDP framework, a proton-type configuration is represented by a signed graph \(G_{\mathrm{p}}\) subject to a collection of structural constraints: three valence cores with prescribed sizes and connectivity patterns, an extensive dynamical relational sea, a fixed total number of internal relations, and the presence of a finite active boundary surface mediating couplings to electron-type \((4,6)\) closures. Even after maximising triangular coherence under these constraints, the associated leakage \(\eta(G_{\mathrm{p}})\) remains strictly positive. The minimal value of \(\eta(G_{\mathrm{p}})\) attainable within the admissible class defines the leakage parameter \(\etaL\), which subsequently reappears in the relational reinterpretation of various coupling constants.
\end{remark}

\subsection{Optimisation functional and admissible proton graphs}

The selection of admissible configurations is governed by a logical optimisation functional that balances coherence, connectivity, and minimality.

\begin{definition}[Optimisation functional]
Let \(\mathcal{G}\) be a collection of finite signed graphs. An \emph{optimisation functional} is a mapping \(\Fopt : \mathcal{G} \to \mathbb{R}\) that depends on the following graph-theoretic quantities:
\begin{itemize}
\item the leakage \(\eta(G)\), which is penalised for larger values;
\item a normalised edge density \(\rho(G) := |E| / |E_{\max}|\), where \(|E_{\max}| = \binom{|V|}{2}\) denotes the number of edges in the complete graph on the vertex set \(V\), thereby rewarding sufficiently high connectivity;
\item a minimality indicator \(m(G) \in \{0,1\}\), which equals \(1\) if \(G\) realises a closed structure that does not rely on any strictly richer superstructure, and \(0\) otherwise.
\end{itemize}
Qualitatively, \(\Fopt(G)\) is assumed to be strictly decreasing as a function of \(\eta(G)\), non-decreasing in \(\rho(G)\) over the relevant range, and larger when \(m(G)=1\). The explicit analytic form of \(\Fopt\) is not required for the purposes of this exposition; we only impose that it favours configurations exhibiting low leakage, sufficiently high connectivity, and minimal closure.
\end{definition}

With this functional in place, one may define the admissible class of proton-type graphs as those that satisfy the prescribed structural constraints and achieve (or approximate) maximal values of the optimisation functional.

\begin{definition}[Admissible proton class and minimal leakage]
Let \(\Gprot \subset \mathcal{G}\) denote the subset of graphs encoding proton-type hierarchical architectures, characterised by the following properties:
\begin{itemize}
\item fixed integer invariants specifying the valence blocks, the relational sea, and the total number of internal relations;
\item prescribed charge and multiplicity constraints (for example, an overall net charge of \(+1\));
\item the presence of a finite active surface that is compatible with coupling to an electron-type \((4,6)\) closure;
\item values of the optimisation functional \(\Fopt\) that are near-maximal under the above constraints.
\end{itemize}
The \emph{minimal leakage parameter} is defined as
\[
 \etaL := \inf_{G \in \Gprot} \eta(G),
\]
and any graph \(G^\ast_{\mathrm{p}} \in \Gprot\) satisfying \(\eta(G^\ast_{\mathrm{p}}) \approx \etaL\) and maximising \(\Fopt\) is selected as a representative proton configuration.
\end{definition}

These preliminaries determine the formal framework in which the subsequent probabilistic and topological constructions are developed. In particular, the leakage functional \(\eta(G)\) is thereby rooted in the combinatorial structure of finite signed graphs, and the proton-level leakage parameter \(\etaL\) is interpreted as a minimal, structurally imposed fraction of non-closure, rather than as an externally adjustable parameter.

\section{Logical Leakage as an Intrinsic Probability}

In the original PDL/LDP framework, the leakage parameter \(\etaL\) is introduced as the minimal closure defect of admissible proton architectures: even after maximising triangular coherence subject to the relevant combinatorial constraints, a strictly positive proportion of triangles necessarily remains violated. In this section we demonstrate that, for any fixed configuration, the same quantity can be reinterpreted as an internal probability of non-closure and that, at the proton level, this probability is not externally prescribed but is instead logically self-sustained by the underlying optimisation principle.

\subsection{Probability Space of Triangles}

We begin by making explicit a natural probability space associated with the set of triangles in a finite signed graph. This construction does not introduce any additional source of randomness; rather, it reformulates the already defined fraction of violated triangles in the language of probability theory.

\begin{definition}[Triangle probability measure]
Let \(G = (V,E,\sigma)\) be a finite signed graph with \(|V| \geq 3\), and let \(\Ttri(G)\) denote the set of all unordered vertex triplets,
\[
 \Ttri(G) := \{ \{i,j,k\} \subset V : i,j,k \text{ pairwise distinct} \}.
\]
We define the sample space by \(\Omega_G := \Ttri(G)\). Since \(\Omega_G\) is finite, we endow it with the full power set as \(\sigma\)-algebra, \(\mathscr{F}_G := 2^{\Omega_G}\). The \emph{triangle probability measure} \(\Prob_G\) is then defined as the uniform probability measure on \(\Omega_G\), given by
\[
 \Prob_G(A) := \frac{|A|}{|\Omega_G|} = \frac{|A|}{\binom{|V|}{3}} \quad \text{for all } A \subseteq \Omega_G.
\]
A random triangle \(T \in \Omega_G\) is said to be \emph{violated} if at least one of the edges between its vertices is absent from \(E\), or if all three edges are present but the product of their associated signs is equal to \(-1\). Otherwise, the triangle is said to be \emph{coherent}.
\end{definition}

Recall that the leakage functional \(\eta(G)\) is originally defined as the normalised count of violated triangles:
\[
 \eta(G) := \frac{\#\{ T \in \Ttri(G) : T \text{ is violated}\}}{|\Ttri(G)|}.
\]
The triangle probability measure permits an immediate probabilistic reinterpretation of this ratio.

\begin{proposition}
For any finite signed graph \(G\) with \(|V| \geq 3\), the leakage \(\eta(G)\) coincides with the probability that a uniformly sampled triangle is violated, that is
\[
 \eta(G) = \Prob_G(T \text{ is violated}),
\]
and, consequently,
\[
 1 - \eta(G) = \Prob_G(T \text{ is coherent}).
\]
\end{proposition}

\begin{proof}
By definition, the set of violated triangles is
\[
 V_G := \{ T \in \Ttri(G) : T \text{ is violated} \},
\]
and the leakage is given by
\[
 \eta(G) = \frac{|V_G|}{|\Ttri(G)|}.
\]
On the other hand, by the construction of the probability measure \(\Prob_G\), we obtain
\[
 \Prob_G(T \text{ is violated}) = \Prob_G(V_G) = \frac{|V_G|}{|\Omega_G|} = \frac{|V_G|}{|\Ttri(G)|}.
\]
Comparing these two expressions yields \(\eta(G) = \Prob_G(T \text{ is violated})\). Since every triangle is either coherent or violated, and these two events are mutually exclusive and collectively exhaustive, it follows that
\[
\Prob_G(T \text{ is coherent}) = 1 - \Prob_G(T \text{ is violated}) = 1 - \eta(G).
\]
This establishes the claim.
\end{proof}

This elementary observation demonstrates that, without modifying any of the underlying combinatorial structure, the leakage functional can be interpreted as the probability of logical failure of closure at the level of elementary triangular cycles.

\subsection{Leakage as a logically self-maintained probability}

We now restrict attention to the class of graphs representing admissible proton architectures. In the PDL/LDP framework, this class is specified by a collection of discrete structural constraints (valence blocks, relational sea, active surface, charge and multiplicity conditions, etc.), together with an optimisation principle that selects configurations exhibiting near-maximal triangular coherence within that class. The leakage parameter \(\etaL\) is then interpreted as the minimal closure defect compatible with these constraints.

\begin{definition}[Admissible proton class and leakage parameter]
Let \(\Gprot\) denote the class of finite signed graphs modelling proton-type hierarchical structures, subject to a fixed family of combinatorial constraints (valence contents, stationary/dynamical partition, existence of an active surface, charge assignment, and related conditions). Let \(\Fopt : \Gprot \to \mathbb{R}\) be an optimisation functional that decreases with the fraction of violated triangles \(\eta(G)\), increases with the normalised edge density, and favours minimally open (i.e.\ maximally closed) structures. We define the \emph{minimal leakage parameter} by
\[
 \etaL := \inf_{G \in \Gprot} \eta(G),
\]
and we select a graph \(G^\ast_{\mathrm{p}} \in \Gprot\) that realises this infimum, or approximates it to within a prescribed tolerance:
\[
 \eta(G^\ast_{\mathrm{p}}) \approx \etaL, \quad \Fopt(G^\ast_{\mathrm{p}}) \ \text{maximal in } \Gprot.
\]
We refer to \(\etaL\) as the \emph{logical leakage} of the proton hierarchy.
\end{definition}

Informally, \(\etaL\) is the smallest achievable fraction of violated triangles among all admissible proton configurations, once all structural constraints have been imposed. Combined with the preceding probabilistic interpretation, this yields an intrinsic probability of non-closure for a ``proton-level'' triangle.

Within the framework, a key conceptual point is that this probability is not a freely adjustable parameter: it is fixed by the same optimisation principle that tends to minimise it. This leads to the following characterisation.

\begin{proposition}[Structural invariance of proton leakage]
Assume that the admissible class \(\Gprot\) and the optimisation functional \(\Fopt\) are fixed once and for all by the underlying axioms and combinatorial constraints of the PDL/LDP framework. Then the minimal leakage \(\etaL\) is a structural invariant of proton-type architectures: it does not depend on any external calibration of a probabilistic parameter, but is instead determined uniquely by the internal logical requirements of closure, coherence optimisation, and minimality.
\end{proposition}

\begin{proof}[Proof (conceptual sketch)]
By construction, the class \(\Gprot\) is specified exclusively in terms of discrete structural data (numbers of vertices and edges in the valence blocks, sizes of the relational sea, total relational content, charge assignments, existence of an active surface, and related quantities) together with qualitative conditions of triangular coherence and connectivity. None of these constraints involves a continuous probabilistic parameter. The optimisation functional \(\Fopt\) is likewise determined by these internal, combinatorial quantities: it penalises violated triangular configurations and rewards connectivity and minimal closed architectures.

Once \(\Gprot\) and \(\Fopt\) are fixed, the value of \(\eta(G)\) for each \(G \in \Gprot\) is completely determined by the combinatorial structure of \(G\). Consequently, the set of achievable leakage values \(\{ \eta(G) : G \in \Gprot \}\) forms a finite (or at most countable) subset of \([0,1]\), and the infimum \(\etaL\) is thereby fixed by the discrete structure of \(\Gprot\). In particular, there is no additional degree of freedom that would allow one to vary \(\etaL\) continuously without altering either the underlying axioms or the admissible architecture itself. Therefore, \(\etaL\) should be interpreted as a structural invariant of proton-type closures: it represents the logically enforced residue of non-closure that persists after coherence has been maximised subject to the given constraints.
\end{proof}

By combining the two propositions, we obtain the central interpretation of this section: at the proton level, the leakage parameter \(\etaL\) coincides with the probability that an elementary triangular cycle, sampled uniformly from within the proton configuration, fails to attain logical closure. Simultaneously, this probability is generated, constrained, and stabilised by the very optimisation principle that is designed to minimise it. Thus, logical leakage does not constitute an externally imposed source of randomness; rather, it is a self-sustained probability of non-closure, dictated by the internal tension between maximal triangular coherence and the finite, hierarchical combinatorial architecture required for the very existence of a proton-type structure.

\section{Topology of Coexistence and Emergent Metric}

Within the PDL/LDP framework, the metric structure is not postulated \textit{a priori}; rather, it is derived from the compatibility relations among closed structures and from the coherence cost associated with sustaining such structures within a common environment. In this section, we formalise this programme by introducing a topology of coexistence on a suitable configuration space and by outlining how a metric notion can be reconstructed as a differential distance induced by variations in coherence cost.

\subsection{Configuration space and compatibility}

We begin by specifying the abstract space in which closed hierarchical structures coexist.

\begin{definition}[Configuration space]
Let \(\Cconf\) denote the set of all closed hierarchical configurations admitted by the PDL/LDP framework. An element \(C \in \Cconf\) may represent, for example, a single elementary \((4,6)\) closure, a proton-like composite structure, a nucleus, or a finite aggregate of such entities. Each configuration is representable by a finite signed graph (or a finite family of such graphs) that satisfies the closure and optimisation constraints at the relevant hierarchical level.
\end{definition}

Informally, two configurations are deemed compatible if they can be embedded into a larger closed structure without inducing an uncontrolled proliferation of coherence violations. This notion can be rigorously characterised in terms of the leakage functional.

\begin{definition}[Logical compatibility]
Let \(C_1, C_2 \in \Cconf\) be two configurations represented by signed graphs \(G_1, G_2\), which may share vertices or be disjoint. Consider a joint configuration \(C_{12}\) obtained by placing \(C_1\) and \(C_2\) in coexistence, together with an appropriate set of inter-configuration edges that satisfy the charge-conservation and structural constraints of the framework. We say that \(C_1\) and \(C_2\) are \emph{logically compatible} if there exists at least one such joint configuration \(C_{12}\) whose associated signed graph \(G_{12}\) satisfies
\[
 \eta(G_{12}) \leq \eta_{\mathrm{max}},
\]
for some fixed threshold \(\eta_{\mathrm{max}} \in (0,1)\) determined by the admissibility criteria of the theory. If no such coexistence configuration exists, we say that \(C_1\) and \(C_2\) are \emph{incompatible}.
\end{definition}

The precise numerical value of \(\eta_{\mathrm{max}}\) is not of primary importance for the present analysis; what is crucial is that the framework specifies an upper bound beyond which a configuration is judged to violate the coherence constraints to such an extent that it can no longer be regarded as physically admissible. Logical compatibility thereby defines a symmetric relation on \(\Cconf\) that specifies which structures can share a common “neighbourhood” without mutually compromising their closure properties.

\begin{definition}[Neighbourhood system]
For each configuration \(C \in \Cconf\), we define its \emph{compatibility neighbourhood} by
\[
 \mathcal{N}(C) := \{ C' \in \Cconf : C' \text{ is logically compatible with } C \}.
\]
We then consider the associated family
\[
 \mathscr{N} := \{ \mathcal{N}(C) : C \in \Cconf \}.
\]
\end{definition}

The family \(\mathscr{N}\) furnishes a natural candidate for a neighbourhood base for a topology on \(\Cconf\). The following proposition outlines the construction of such a topology.

\begin{proposition}[Coexistence topology]
Suppose that logical compatibility is reflexive (every configuration is compatible with itself) and closed under finite conjunctions, in the sense that if \(C_1, C_2, C_3 \in \Cconf\) are such that any two among them are pairwise compatible, then there exists a configuration \(C_{123}\) in which all three coexist with \(\eta(G_{123}) \leq \eta_{\mathrm{max}}\). Under these assumptions, the family \(\mathscr{N}\) generates a topology \(\mathscr{T}\) on \(\Cconf\), called the \emph{coexistence topology}, for which each \(\mathcal{N}(C)\) is an open neighbourhood of \(C\).
\end{proposition}

\begin{proof}[Proof (sketch)]
Define a collection \(\mathscr{T}\) of subsets of \(\Cconf\) by declaring \(U \subseteq \Cconf\) to be open if, for every \(C \in U\), there exists a configuration \(C' \in \Cconf\) such that \(C \in \mathcal{N}(C') \subseteq U\). By construction, both \(\emptyset\) and \(\Cconf\) satisfy this condition. Furthermore, the union of any family of sets with this property again satisfies the same condition, and therefore \(\mathscr{T}\) is closed under arbitrary unions.

To verify stability under finite intersections, let \(U_1, U_2 \in \mathscr{T}\) and \(C \in U_1 \cap U_2\). By definition of \(\mathscr{T}\), there exist configurations \(C_1, C_2 \in \Cconf\) such that
\[
C \in \mathcal{N}(C_1) \subseteq U_1, \qquad C \in \mathcal{N}(C_2) \subseteq U_2.
\]
By reflexivity and symmetry of logical compatibility, the configurations \(C, C_1, C_2\) are pairwise compatible. By the finite conjunction assumption, there exists a configuration \(C_{12}\) in which \(C, C_1, C_2\) coexist with admissible leakage, i.e., with \(\eta(G_{12}) \leq \eta_{\mathrm{max}}\). Consequently, \(\mathcal{N}(C_{12})\) is a neighbourhood of \(C\) contained in \(U_1 \cap U_2\), which implies that \(U_1 \cap U_2 \in \mathscr{T}\). Thus, \(\mathscr{T}\) satisfies the axioms of a topology on \(\Cconf\).

Finally, by construction, each \(\mathcal{N}(C)\) is open and contains the configuration \(C\), and therefore serves as a neighbourhood of \(C\) in the topological space \((\Cconf,\mathscr{T})\).
\end{proof}

The induced topological structure thereby captures a purely relational notion of “coexistence without excessively severe violations of coherence.” Crucially, no metric or geometric structure is presupposed at this stage: the topology arises solely from the underlying logical compatibility conditions.

\subsection{Metric as coherence-cost distance}

Once a topology of coexistence has been specified, it becomes natural to introduce a quantitative notion of separation between configurations, not in terms of predetermined spatial coordinates, but rather in terms of the additional coherence cost required to sustain them within a common neighbourhood. This consideration leads to the construction of a pseudometric derived from the leakage functional.

\begin{definition}[Coherence-cost distance]
Let \(C_1, C_2 \in \Cconf\) be two configurations, represented by graphs \(G_1\) and \(G_2\), respectively. Consider the set \(\mathcal{J}(C_1,C_2)\) of all admissible joint configurations \(C_{12}\) in which \(C_1\) and \(C_2\) coexist, and whose signed graphs \(G_{12}\) extend \(G_1\) and \(G_2\) while satisfying the structural constraints imposed by the framework. The \emph{coherence-cost distance} between \(C_1\) and \(C_2\) is defined by
\[
 d(C_1,C_2) := \inf_{C_{12} \in \mathcal{J}(C_1,C_2)} \bigl( \eta(G_{12}) - \max\{\eta(G_1), \eta(G_2)\} \bigr)_+,
\]
where \((x)_+ := \max\{x,0\}\) denotes the positive part. If no admissible joint configuration exists, we set \(d(C_1,C_2) := +\infty\).
\end{definition}

Informally, \(d(C_1,C_2)\) quantifies the minimal additional leakage that the relational framework must accommodate when \(C_1\) and \(C_2\) are required to coexist, relative to the worst leakage incurred by each configuration in isolation. If the two configurations can coexist without increasing the leakage beyond \(\max\{\eta(G_1),\eta(G_2)\}\), then \(d(C_1,C_2) = 0\); otherwise, the distance represents the smallest unavoidable overhead in terms of coherence cost.

\begin{proposition}[Pseudometric properties]
Under mild regularity assumptions on the class of admissible joint configurations (specifically, closure under relabelling and the existence of sufficiently symmetric embeddings), the function \(d : \Cconf \times \Cconf \to [0,+\infty]\) defined above is a pseudometric on the subset of \(\Cconf\) where \(d\) is finite. In particular, it satisfies:
\begin{enumerate}[label=(\roman*)]
\item \(d(C_1,C_2) \geq 0\) and \(d(C_1,C_1) = 0\) for all \(C_1, C_2\),
\item \(d(C_1,C_2) = d(C_2,C_1)\) for all \(C_1, C_2\),
\item \(d(C_1,C_3) \leq d(C_1,C_2) + d(C_2,C_3)\) whenever all three distances are finite.
\end{enumerate}
\end{proposition}

\begin{proof}[Proof (heuristic outline)]
Non-negativity and reflexivity follow immediately from the definition: taking the positive part ensures \(d \geq 0\), and choosing \(C_{11} = C_1\) as a degenerate joint configuration yields \(d(C_1,C_1) = 0\). Symmetry holds because the definition of \(\mathcal{J}(C_1,C_2)\) is symmetric in \(C_1\) and \(C_2\), and the quantity \(\eta(G_{12}) - \max\{\eta(G_1),\eta(G_2)\}\) is invariant under exchanging \(G_1\) and \(G_2\).

For the triangle inequality, consider three configurations \(C_1, C_2, C_3\) such that the pairwise distances are finite. By the definition of the infimum, for every \(\varepsilon > 0\) there exist joint configurations \(C_{12} \in \mathcal{J}(C_1,C_2)\) and \(C_{23} \in \mathcal{J}(C_2,C_3)\) whose leakages approximate the corresponding infima up to \(\varepsilon\). Under mild composability assumptions (namely, that these joint configurations can be embedded into a larger joint configuration \(C_{123}\) without incurring uncontrolled additional violations), the leakage of the associated graph \(G_{123}\) can be bounded in terms of the leakages of \(G_{12}\), \(G_{23}\), and the individual graphs. This provides an upper bound on
\[
\eta(G_{123}) - \max\{\eta(G_1),\eta(G_3)\}
\]
that is essentially the sum of the two intermediate coherence penalties, up to arbitrarily small corrections controlled by \(\varepsilon\). Passing to infima then yields
\[
d(C_1,C_3) \leq d(C_1,C_2) + d(C_2,C_3).
\]

A fully rigorous proof would require a precise specification of the class of admissible joint embeddings and their structural properties, which lies beyond the scope of this note. Conceptually, however, the argument parallels standard constructions of pseudometrics from cost functionals in combinatorial optimisation.
\end{proof}

The coherence-cost distance thereby furnishes a quantitative refinement of the coexistence topology. Small values of \(d(C_1,C_2)\) indicate that the two configurations can be sustained within the same region of the relational network without necessitating substantial additional leakage, whereas large values correspond to a pronounced deformation of the coherence budget. In the limit in which many configurations coexist and the network of compatibility relations becomes dense, the function \(d\) can be interpreted as the discrete precursor of an emergent metric. From this vantage point, what is conventionally described as “spatial distance” is reinterpreted as the minimal coherence-cost differential required to support the coexistence of structures, and gravitational phenomena are mapped to large-scale variations of this cost across the configuration space.

\section{A Logical Entropic Paradox}

The probabilistic and topological reformulations developed above enable a comparison between two distinct notions of entropy: a local entropy of closure, associated with the distribution of coherent and violated triangles within a single configuration, and a global or topological entropy, associated with the organisation of coexistence in the emergent configuration space. In this section, we demonstrate, at a schematic level, how an optimisation principle that drives the minimisation of local entropy can, at the same time, induce an increase in global entropy as additional structures are aggregated. We refer to this tension as a logical entropic paradox.

\subsection{Local entropy of closure versus global entropy of coexistence}

We begin by introducing a minimal notion of local logical entropy, formulated in terms of the leakage probability associated with triangular substructures.

\begin{definition}[Local logical entropy]
Let \(G\) be a finite signed graph representing a closed configuration, and let \(\eta(G) \in [0,1]\) denote its leakage, interpreted as the probability that a uniformly randomly sampled triangle in \(G\) is violated. The \emph{local logical entropy} of \(G\) is defined as the binary Shannon entropy
\[
 S_{\mathrm{loc}}(G) := H(\eta(G)) := - \eta(G) \log \eta(G) - (1 - \eta(G)) \log (1 - \eta(G)),
\]
with the standard convention \(0 \log 0 := 0\). More generally, this definition can be refined by considering the full probability distribution over the possible types of triangles (for instance, coherent triangles, triangles violated by missing edges, triangles violated by sign frustration, and so forth) and computing the Shannon entropy of that distribution. For the purposes of the present work, the binary approximation is sufficient.
\end{definition}

The local entropy \(S_{\mathrm{loc}}(G)\) quantifies the uncertainty associated with the closure status of a randomly selected triangle within a given configuration. For a perfectly closed structure, such as the elementary \((4,6)\) block, we have \(\eta(G) = 0\) and consequently \(S_{\mathrm{loc}}(G) = 0\); in this case, no uncertainty is present because every triangle is coherent. For realistic composite structures, such as the proton, one typically finds \(0 < \etaL \ll 1\), implying that \(S_{\mathrm{loc}}(G)\) is small but strictly positive. This captures the fact that an overwhelming majority of triangles are coherent, while a nonzero fraction fails to close.

To formulate a notion of global entropy, we analyse how leakage is distributed over regions of configuration space when multiple structures coexist.

\begin{definition}[Global topological entropy (schematic)]
Let \((\Cconf,\mathscr{T})\) denote the coexistence topological space, and let \(\mu\) be a probability measure on \(\Cconf\) that characterises the statistical prevalence of different configurations in a given physical regime (for example, within a macroscopic sample of ordinary matter or within a specified cosmological model). The leakage functional \(\eta : \Cconf \to [0,1]\) induces a random variable \(\eta(C)\) on the probability space \((\Cconf,\mu)\). The \emph{global topological entropy} of the regime is then defined, in a minimal form, as
\[
 S_{\mathrm{glob}} := H(\eta(C)) := - \int_0^1 p(\lambda) \log p(\lambda) \, d\lambda,
\]
where \(p\) denotes the probability density (or, in the discrete case, the probability mass function) of the leakage values induced by \(\mu\). More refined definitions may incorporate spatial or configurational correlations of leakage via the coherence-cost distance \(d\) or via suitable partitions of \(\Cconf\). In the present discussion, however, we interpret \(S_{\mathrm{glob}}\) as a measure of the heterogeneity of leakage levels realised across the ensemble of coexisting configurations.
\end{definition}

Intuitively, \(S_{\mathrm{glob}}\) attains low values when nearly all configurations realised within a given regime exhibit the same leakage value (for instance, when a single structural type is present), and it increases as the distribution of leakage values becomes more heterogeneous. Consequently, \(S_{\mathrm{glob}}\) is sensitive to the richness of the hierarchical and compositional organisation of matter and, indirectly, to the complexity of the emergent metric induced by the cumulative coherence cost.

Taken together, the two entropy notions highlight a structural tension: the optimisation functional tends to minimise \(\eta(G)\) for each configuration, thereby driving \(S_{\mathrm{loc}}(G)\) towards low values, whereas the proliferation of distinct configurations and their coexistence in a shared topological space tends to increase the diversity of leakage patterns and thus to raise \(S_{\mathrm{glob}}\). The following proposition illustrates this phenomenon in a simplified, schematic setting.

\begin{proposition}[Schematic entropic paradox]
Consider a family \(\{G_i\}_{i=1}^N\) of admissible configurations in \(\Cconf\), each locally optimised so that \(\eta(G_i) \approx \eta_{\ast}\) for some small leakage value \(\eta_{\ast} \in (0,1)\). Assume that:
\begin{enumerate}[label=(\alph*)]
\item at the level of individual configurations, the optimisation functional \(\Fopt\) favours smaller values of \(\eta(G)\), so that any perturbation that increases \(\eta(G_i)\) is dynamically disfavoured;
\item when configurations are aggregated and forced to coexist, the resulting joint configurations explore a family of effective leakage values \(\tilde{\eta}\) that deviate from \(\eta_{\ast}\) in a manner that depends on the local density and hierarchical organisation of the \(G_i\);
\item the statistical distribution of these effective leakage values, regarded as a random variable over the aggregate, broadens as \(N\) increases and as the organisation of the aggregate becomes more heterogeneous.
\end{enumerate}
Then, as the number and structural diversity of coexisting configurations increase, the average local entropy \(\frac{1}{N}\sum_i S_{\mathrm{loc}}(G_i)\) remains small (since each \(\eta(G_i) \approx \eta_{\ast}\)), whereas the global topological entropy \(S_{\mathrm{glob}}\), associated with the distribution of effective leakage values, increases. In other words, local coherence optimisation coexists with a growth of global entropic complexity in the configuration space.
\end{proposition}

\begin{proof}[Proof (informal illustration)]
By assumption, each \(G_i\) individually satisfies \(\eta(G_i) \approx \eta_{\ast}\). For fixed \(\eta_{\ast}\), the corresponding local entropy, defined as \(S_{\mathrm{loc}}(G_i) = H(\eta(G_i))\), is therefore approximately constant and takes a relatively small value. Replicating such configurations without modifying their internal structure leaves this average essentially unchanged.

When these configurations are assembled into a larger composite system and required to coexist under additional constraints (for instance, within nuclei, bulk matter, or astrophysical objects), the effective leakage experienced by different regions of the composite system need not remain uniform. Spatial variations in hierarchical packing, heterogeneities in nuclear composition, or differences in surface-to-volume ratios can induce small shifts in the local coherence requirements needed to maintain global compatibility. This induces a family \(\{\tilde{\eta}_\alpha\}\) of effective leakage parameters associated with distinct coarse-grained regions or emergent substructures.

If the internal organisation of the aggregate is trivial—such as in an idealised homogeneous crystal composed of identical units—this spectrum may remain narrowly concentrated around \(\eta_{\ast}\), in which case the global entropy \(S_{\mathrm{glob}}\) remains small. However, as additional levels of structural hierarchy are introduced (protons forming nuclei, nuclei forming atoms and molecules, molecules forming complex solids and fluids, and these in turn forming gravitationally bound macroscopic systems), the interaction between local closure constraints and global compatibility conditions generates an increasingly rich set of effective leakage values. Consequently, the distribution \(p(\lambda)\) of leakage levels across the aggregate broadens, and the associated Shannon entropy \(S_{\mathrm{glob}} = H(\eta(C))\) increases.

Throughout this hierarchical construction, the underlying optimisation principle continues to drive \(\eta(G)\) towards its minimum at each local structural scale, so that each fundamental building block remains close to the reference leakage \(\eta_{\ast}\) and retains low local entropy. The growth of \(S_{\mathrm{glob}}\) therefore does not result from any explicit ``randomisation'' of the individual structures; rather, it arises from the combinatorial multiplicity of ways in which numerous nearly optimal closures can be arranged within a common coexistence domain. In this precise sense a logical entropic paradox appears: the same tendency towards minimal local leakage entails, once a sufficiently large number of structures are allowed to coexist, a proliferation of distinct macroscopic patterns by which a small but irreducible degree of non-closure can be distributed over large scales.
\end{proof}

This apparent paradox does not indicate any inconsistency within the theoretical framework; rather, it exemplifies a well-known phenomenon in purely relational formalisms: the minimization of an underlying cost functional at the microscopic scale is fully compatible with, and in fact typically gives rise to, a high degree of structural and entropic complexity at macroscopic scales. In the context considered here, the relevant cost is the leakage-based coherence functional, the microscopic constituents are closed hierarchical structures such as protons and nuclei, and the macroscopic complexity emerges both in the induced effective metric and in the statistical distribution of coherence costs over the configuration space.

\section{Connections to Physical Quantities}

The probabilistic and topological reformulations developed above are designed to remain strictly internal to the logical framework, without presupposing any specific physical interpretation. Nonetheless, they acquire concrete significance once they are mapped back onto the quantities that PDL/LDP already recast in relational terms, such as Born-type probabilities, gravitational coupling constants, and thermal granularity scales. In this section, we outline these connections—without attempting a complete rederivation of the numerical results—and emphasise how the concept of logically self-maintained leakage furnishes a unifying interpretative framework.

In the PDL/LDP framework, the complete signed graph on four vertices and six edges, denoted \((4,6)\), is singled out as the first minimal elementary stationary structure. More precisely, there exists an assignment of edge signs such that all triangular subgraphs are coherent, and this coherence is preserved under a global binary pulsation of the edge signs. Within this scheme, the \((4,6)\) structure is proposed as a logical prototype of an electron at rest. In the terminology introduced above, the electron-type \((4,6)\) closure satisfies \(\eta(G_{4,6}) = 0\), implying that both the local leakage probability and the local entropy are identically zero: every triangle closes without defect, and the internal reference frame is completely self-contained.

When an electron-type block couples to a proton-type hierarchical structure, the interaction is mediated through mixed triangles whose vertices lie partly in the \((4,6)\) block and partly on the proton’s active surface. In the combinatorial representations employed in PDL/LDP, these mixed triangles carry a flux of coherence that, upon coarse-graining, can be interpreted as probability amplitudes for locating the electron in various regions of an emergent configuration space. Within this perspective, the following representation is natural: internally, the \((4,6)\) block furnishes a strictly closed reference structure, exhibiting no intrinsic leakage, whereas externally, the pattern of coherent and violated mixed triangles encodes the manner in which a small portion of the proton’s logically self-sustained leakage is captured, redirected, and redistributed along distinct channels.

From a formal perspective, one may consider the following minimal framework. Let \(G_{\mathrm{p}}\) denote a proton-type graph characterised by a leakage parameter \(\etaL\), and let \(G_{e}\) denote a \((4,6)\) graph with vanishing leakage, \(\eta(G_{e}) = 0\). Consider a family of joint configurations \(\{G_{e\mathrm{p}}(x)\}\), indexed by a discrete parameter \(x\) that labels the admissible “positions’’ or modes of the electron relative to the proton. Each such joint configuration induces a specific distribution of mixed triangles and thereby determines a configuration-dependent leakage \(\eta(G_{e\mathrm{p}}(x))\). This leakage can be interpreted as quantifying how the intrinsic leakage of the proton is partitioned among the distinct patterns of mixed triangles.

By normalising these contributions over all values of \(x\), one is naturally led to define a probability measure \(\Prob(x)\) on the discrete configuration space. In suitable limiting regimes, this measure is expected to reproduce \(|\psi(x)|^2\) for simple bound states. Within this interpretation, Born-type probabilities no longer require an additional probabilistic axiom; instead, they arise as a particular projection of the logically self-maintained leakage of composite structures onto the space of mixed-triangle couplings with elementary closures.

A full realisation of this programme would necessitate a detailed combinatorial analysis of the joint dynamics of \((4,6)\) graphs and proton-type graphs, together with a systematic comparison to the standard Schrödinger framework. The central claim advanced here is more limited in scope: once leakage is understood as an intrinsic probability of non-closure, it becomes natural to interpret the statistical weights in quantum-like descriptions as expressions of how this probability is redistributed across mixed relational channels, rather than as manifestations of external stochasticity imposed on an otherwise deterministic logical substrate.

\subsection{Gravitational coupling and hierarchical filtering}

A central assertion of the PDL/LDP framework is that an effective gravitational coupling constant can be formulated in terms of the same relational quantities that characterise nucleonic structures, specifically the proton mass, the fine-structure constant, and the leakage parameter. More precisely, an expression of the form
\[
 G_{\mathrm{eff}} \sim \frac{\hbar c}{m_p^2} \, \alpha \, \etaL^{18}
\]
is proposed, modulo numerical conventions and higher-order refinements. Here \(m_p\) denotes the proton mass, \(\alpha\) is the fine-structure constant reinterpreted as a surface-to-volume coherence fraction, and the exponent \(18\) encodes the number of hierarchical coherence filters that a local leakage event must traverse before it manifests as a universal, large-scale metric deformation.

Within the probabilistic formalism introduced earlier, \(\etaL\) is the logical probability that a proton-level triangular configuration fails to close, and \(\etaL^{18}\) may be interpreted as an effective multi-level probability that a coherence packet escapes closure not only at the proton scale, but throughout the entire sequence of filters linking proton-scale defects to macroscopic gravitational phenomena. The specific value of the exponent reflects the number of distinct, structurally independent constraints that attenuate and redistribute leakage as one ascends the hierarchy (from \((4,6)\) bricks to protons, from protons to nuclei, and from nuclei to ordinary matter and the emergent metric). Each such filter partially absorbs or reconfigures the leakage into shorter-range interactions or internal stabilisation mechanisms, leaving only a residual component capable of contributing to long-range, universal curvature.

From this perspective, the extreme weakness of gravity relative to the other fundamental interactions ceases to be puzzling: it is a direct consequence of the very low probability that a coherence defect survives all filtering stages and is ultimately projected into the emergent metric. The gravitational coupling constant is thereby reinterpreted as a measure of the frequency of such residual leakage events per unit of action, mass, and surface coherence, rather than as an independent, primitive parameter. The probabilistic/topological reformulation thus renders more precise the heuristic statement that ``gravity is the residue of closure failure'': gravity is construed as the macroscopic imprint of a logically self-maintained probability that persists after coherence has been maximised and filtered across all accessible scales.

\subsection{Thermal granularity and Boltzmann's constant}

The same leakage parameter that appears in the expression for the effective gravitational coupling also enters into proposed expressions for Boltzmann’s constant \(k_B\) within the PDL/LDP framework. In this context, \(k_B\) is reinterpreted as quantifying the minimal granularity of coherence associated with collective excitations in ensembles of hierarchical structures, rather than serving solely as a phenomenological conversion factor between temperature and energy.

In the terminology employed here, this reinterpretation can be formulated as follows. At the level of a single proton, \(\etaL\) characterizes the probability that a relational triangle fails to close and, consequently, that a fraction of the internal coherence budget is effectively projected into the relational environment. When many such structures constitute an ensemble (for example, in a gas or a solid), these rare leakage events function as microscopic “kicks’’ that enable coherence to be redistributed among different modes and degrees of freedom. The average energy associated with these redistributions, per degree of freedom and per unit excitation probability, establishes an energy scale that can be identified with \(k_B T\) in equilibrium regimes.

Within this perspective, temperature quantifies the average occupation of coherence cycles that have been driven beyond their minimal stationary regime, while Boltzmann’s constant specifies the conversion factor between the logical granularity of these excitations and the macroscopic energy scale. The fact that the candidate expression for \(k_B\) can be written exclusively in terms of the relational integers specifying the electron and proton architectures, together with the same leakage parameter \(\etaL\) that determines gravitational coupling, suggests that thermal granularity and gravitation may share a common logical origin in the propagation and accumulation of coherence defects across hierarchical levels.

Collectively, these connections to Born-type probabilities, gravitational coupling, and thermal granularity highlight the central role of logically self-maintained leakage. When reformulated as an intrinsic probability of non-closure, this leakage furnishes a unifying conceptual mechanism that links quantum statistics, emergent metric curvature, and thermodynamic discreteness within a single relational substrate.

\section{Discussion}

The constructions presented in this note do not alter any of the combinatorial constituents of the PDL/LDP framework. Rather, they render explicit a probabilistic and topological layer that was only implicit in the original formulation: the leakage functional \(\eta\) is reformulated as an internal probability of non-closure over triangles, and the emergent metric is reinterpreted as a coherence-cost distance on a space of mutually coexisting closures. In this section, we summarise the conceptual and structural benefits afforded by this reformulation and outline several directions for further theoretical development.

\subsection{Conceptual gains}

The first conceptual gain consists in recognising that leakage is not merely a perturbative correction to an otherwise perfectly closed hierarchy, but rather a well-defined logical probability associated with elementary cycles. By endowing the set of triangles of a finite signed graph with a natural uniform measure, the proportion of violated triangles can be interpreted as the probability that an internal reference loop fails to close. This interpretation does not require any supplementary source of randomness; it is entirely determined by the underlying combinatorial structure together with the coherence assignment.

The second gain pertains to the status of the leakage parameter \(\etaL\) at the proton level. In the original PDL/LDP framework, \(\etaL\) arises as a minimal closure defect consistent with a set of discrete structural constraints and with a given optimisation functional. Once recast in probabilistic terms, it becomes evident that \(\etaL\) exhibits the behaviour of a self-maintained probability: it is both generated and stabilised by the very principle that tends to minimise it, and its numerical value is fixed by the internal architecture rather than through an external calibration of a stochastic parameter. This supports the interpretation of probability not as an ad hoc supplement to the logical framework, but as an emergent property of imperfect closure in finite relational structures.

The third gain lies in the elucidation of the emergent metric from a topological perspective. Instead of treating distance as a primitive, independent notion, the coexistence topology together with the coherence-cost pseudometric yields a purely relational characterisation: two configurations are proximate if they can share a neighbourhood without incurring a substantial increase in leakage. The resulting metric structure can thus be viewed as a coarse-grained representation of the distribution of coherence costs over configuration space, and gravitational phenomena appear as large-scale manifestations of gradients in this cost, rather than as intrinsic attributes of a pre-specified geometric substratum.

\subsection{Logical entropy and large-scale complexity}

The notion of a logical entropic paradox, introduced here in schematic form, also clarifies the mechanism by which complexity arises within the PDL/LDP framework. At the local level, each closed structure minimises its leakage and thereby its local entropy: an elementary \((4,6)\) block exhibits zero leakage, while empirically relevant composite structures such as protons or nuclei possess very small but non-zero leakage. When a large number of such structures are permitted to coexist and to organise into increasingly intricate hierarchical arrangements, the distribution of effective leakage values over the accessible configuration space widens. This broadening induces an increase in global, or topological, entropy, despite the fact that each constituent remains individually optimised.

An analogous pattern is well known from statistical mechanics and information theory: microscopically ordered or low-entropy units can collectively generate macroscopically high-entropy ensembles when they can be combined in a large number of distinct configurations. The present framework extends this intuition by providing a relational, logically grounded account of the phenomenon, in which both the microscopic cost and the macroscopic complexity are formulated in terms of a common set of coherence and closure constraints on signed graphs.

\subsection{Limitations and open problems}

Several significant limitations of the present work should be explicitly acknowledged.  

First, the probabilistic and entropic constructions remain at a schematic level. The triangle probability measure is mathematically well-defined, but the global topological entropy is, at this stage, only outlined as a functional of the leakage distribution over the configuration space. A more complete and rigorous treatment would require explicit and well-justified choices of measures on \(\Cconf\) and of coarse-grainings that are compatible with the emergent metric structure, together with concrete computations in suitably chosen toy models.

Second, the coherence-cost pseudometric has been introduced in a largely heuristic manner. Although its basic properties as a pseudometric can be plausibly motivated by analogy with cost-based distances in combinatorial optimisation, a fully rigorous formulation would require a precise specification of the class of admissible joint embeddings of configurations and a detailed account of the behaviour of leakage under composition. Progress in this direction is closely tied to obtaining a more explicit characterisation of the admissible proton, nuclear, and matter graphs within the PDL/LDP framework.

Third, while the connections to Born-type probabilities, gravitational coupling, and thermal granularity have been sketched at a conceptual level, they have not yet been developed into complete derivations. In particular, deriving the Born rule from fluxes of mixed triangles in the \((4,6)\)-proton system, obtaining the exponent \(18\) in the gravitational coupling from a purely combinatorial enumeration of hierarchical filters, and determining the numerical value of \(\etaL\) from first principles rather than via empirical calibration all remain open problems. Resolving these issues would advance the current reformulation from a largely conceptual clarification to a quantitatively predictive and testable framework.

The present analysis is, in effect, restricted to a static regime: it concentrates on stationary closures and their coexistence, without systematically incorporating dynamical features such as constraint events, genuinely non-equilibrium processes, or the temporal formation and dissolution of closures. A natural extension would be to generalize the probabilistic and topological framework developed here so as to encompass these intrinsically dynamical phenomena, which may in turn provide novel insights into cosmological scenarios and into the emergence of temporal asymmetry within a fundamentally relational framework.

Taken as a whole, the discussion indicates that the probabilistic and topological reformulations advanced in this work are not merely an optional reinterpretation, but rather constitute a natural completion of the original PDL/LDP programme. They render explicit the structural interplay among coherence, closure, leakage, and probability, and they delineate specific mathematical problems whose resolution would advance the framework towards a more fully articulated relational theory of microphysics, gravitation, and thermodynamics.

\section{Conclusion}

The objective of this note has been to render explicit a probabilistic and topological layer that is already implicit in the PDL/LDP framework. Beginning with the established definition of the leakage functional \(\eta(G)\) as the normalised count of violated triangles in a finite signed graph, we have shown that it admits a reinterpretation as an intrinsic probability of non-closure for elementary triangular cycles. In particular, the proton-level leakage parameter \(\etaL\) arises as a logically self-sustaining probability: it is determined endogenously by the same coherence-optimisation principle that seeks to minimise it, rather than being imposed by an external stochastic postulate.

Building on this reinterpretation, we have introduced a coexistence topology on the space of closed hierarchical configurations, generated by logical compatibility relations, and we have outlined how an emergent metric structure can be reconstructed as a coherence-cost pseudometric. In this framework, distance no longer appears as a primitive geometric datum, but instead as a quantitative measure of the additional leakage required to sustain multiple closures within a common relational environment. Gravitation can, in turn, be understood as a large-scale modulation of this coherence cost, governed by the extremely small probability that proton-level leakage persists through a finite chain of hierarchical filters.

We have also contrasted a notion of local logical entropy, defined in terms of the leakage probability within individual configurations, with a global or topological entropy, defined in terms of the distribution of effective leakage levels across the overall configuration space. This contrast gives rise to a logical entropic paradox: the very optimisation procedures that tend to minimise local leakage and thereby reduce local entropy for each individual structure simultaneously allow for, and in fact generically induce, an increase in global entropic complexity when many such structures coexist and self-organise into intricate hierarchical arrangements.

Furthermore, we have outlined how this probabilistic and topological reformulation interfaces with physical quantities that the PDL/LDP framework already recasts in relational terms: Born-type probabilities as redistributions of leakage across mixed-triangle channels; gravitational coupling as a multi-level probability of residual leakage encoded in an exponent such as \(18\); and Boltzmann’s constant as a measure of thermal granularity associated with coherence defects in ensembles of nucleonic structures. These connections are, at this stage, schematic, but they indicate that quantum statistics, gravitation, and thermodynamics may be interpreted as distinct large-scale manifestations of a common, logically mandated imperfection of closure in a finite relational network.

From this perspective, the present contribution is best regarded as a conceptual completion rather than a revision of the PDL/LDP programme. It elucidates how coherence, closure, leakage, probability, topology, and metric structure can be integrated into a unified logical framework, and it delineates specific mathematical and phenomenological problems whose resolution would constitute a significant advance toward a fully relational description of physical reality.
\cite{Laubscher2026_PDL_long}

\bibliographystyle{plain}
\bibliography{pdl_leakage_topology}

\end{document}
